%% ============================================================
%%  CHAPTER 10 — RESULTS AND ANALYSIS
%% ============================================================
\chapter{Results and Analysis}
\label{ch:results}

This chapter presents the detailed evaluation results of the final
production model (the \acrshort{dtw} hybrid architecture,
Chapter~9), all under the honest Leave-One-Session-Out
(\acrshort{loso}) protocol established in that chapter.
Section~\ref{sec:res_headline} restates the headline result and
reconciles the two ways of aggregating the LOSO folds.
Section~\ref{sec:res_perfold} breaks that result down fold by
fold. Section~\ref{sec:res_perclass} presents the per-class
performance breakdown. Section~\ref{sec:res_confusion} presents
and analyses the confusion matrix. Section~\ref{sec:res_ablation}
presents the complete experiment log as a full ablation table,
consolidating every architectural and feature decision made across
Chapters~6--9 into a single reference. Section~\ref{sec:res_realtime}
summarises the real-time operating characteristics of the deployed
model. Section~\ref{sec:res_summary} summarises the chapter.

%% ----------------------------------------------------------
\section{Headline Result}
\label{sec:res_headline}
%% ----------------------------------------------------------

The final production model --- the three-way hybrid architecture
of Chapter~9, combining hand-crafted-feature routing (Stage~1),
\acrshort{dtw} shape-matching for blink sub-classification
(Stage~2a), and Riemannian covariance geometry for muscle
sub-classification (Stage~2b) --- achieves:

\begin{center}
\Large \textbf{91.3\% accuracy, macro-F1 0.913}\\
\normalsize (unweighted mean across the 3 LOSO folds, as reported
during architecture development, Chapter~9)\\[4pt]
\Large \textbf{91.8\% accuracy, macro-F1 0.918}\\
\normalsize (pooled across all 1{,}246 held-out predictions,
Section~\ref{sec:res_perclass})
\end{center}

Both figures describe the same trained model and the same
\acrshort{loso} protocol; they differ only in how the three
per-session folds are aggregated. The 91.3\% figure
(Chapter~9) is the \textbf{unweighted mean} of the three per-fold
accuracies; the 91.8\% figure is computed by \textbf{pooling} every
held-out prediction from all three folds into one set before
scoring. The two are consistent to within 0.5 points --- the pooled
figure is very slightly higher because Session3, the largest fold
(509 of 1{,}246 epochs) and the model's best-performing fold
(95.7\%, Section~\ref{sec:res_perfold}), is weighted by its epoch
count under pooling rather than counted equally with the two
smaller sessions. Table~\ref{tab:perfold_accuracy} in
Section~\ref{sec:res_perfold} gives the exact per-fold numbers
behind both figures. Since the pooled evaluation was run directly
against the production \code{DtwRiemannHybridClassifier} used by
\code{train\_dtw\_hybrid.py} (rather than summarised from prior
documentation), it is treated as the authoritative source for the
detailed per-class and confusion-matrix results in this chapter;
Chapter~9 retains the 91.3\% figure as originally reported during
development, for historical continuity with the progression
narrative told there.

This is, per the reference table established in Chapter~9
(Table~\ref{tab:honest_numbers}), the family of numbers that should
be cited as \emph{the} accuracy of the system: computed under the
strictest available protocol (an entire unseen recording session,
cross-day, with no epoch-overlap or preprocessing leakage), and
answering the deployment-relevant question --- how well the system
performs on a new day without recalibration.

%% ----------------------------------------------------------
\section{Per-Fold Breakdown}
\label{sec:res_perfold}
%% ----------------------------------------------------------

Table~\ref{tab:perfold_accuracy} gives the accuracy of each
architecture (Chapter~9) on each of the three \acrshort{loso}
folds individually, computed by \code{eval\_loso.py} against the
same session splits used throughout Chapter~9.

\begin{table}[htbp]
\centering
\caption{Per-fold LOSO accuracy by architecture and held-out
session. Each column is one fold: train on the other two sessions,
test on the named session. ``dtw-hybrid'' is the final production
model.}
\label{tab:perfold_accuracy}
\renewcommand{\arraystretch}{1.3}
\begin{tabular}{l c c c c}
\toprule
\textbf{Held-out session} & \textbf{Gated epochs} &
\textbf{Flat RF} & \textbf{Hierarchical} & \textbf{DTW-hybrid} \\
\midrule
Session1 & 368 & 84.8\% & 85.3\% & 90.2\% \\
Session2 & 369 & 83.5\% & 82.7\% & 88.1\% \\
Session3 & 509 & 87.0\% & 87.2\% & 95.7\% \\
\midrule
\textbf{Unweighted mean} & --- & 85.1\% & 85.1\% &
\textbf{91.3\%} \\
\bottomrule
\end{tabular}
\end{table}

Three things are worth noting in this table:

\begin{itemize}[leftmargin=*]
    \item The unweighted-mean row reproduces exactly the headline
    figures reported throughout Chapter~9 (85.1\% flat, 85.1\%
    hierarchical, 91.3\% DTW-hybrid), confirming the per-fold
    numbers behind those figures.
    \item \textbf{DTW-hybrid outperforms both simpler architectures
    on every single fold}, not just on average --- the improvement
    from hierarchical routing plus Riemannian/DTW branches is
    consistent across all three recording sessions, not an artefact
    of one favourable session.
    \item \textbf{Session3 is both the largest and the
    best-performing fold} for every architecture (509 gated
    epochs, 95.7\% for DTW-hybrid), which is why the pooled
    accuracy in Section~\ref{sec:res_headline} (91.8\%) is slightly
    higher than the unweighted per-fold mean (91.3\%): pooling
    naturally gives Session3's predictions more influence over the
    final score, in proportion to how much data it contributes.
\end{itemize}

\textit{Methodological note: producing this table required a small
correction to the development-time evaluation script. }
\code{eval\_loso.py}\textit{ originally only fit and compared the
flat and hierarchical models; the DTW-hybrid numbers previously
reported (Chapter~9) came from a separate one-off script
(}\code{eval\_dtw\_blink.run\_end\_to\_end()}\textit{) that
reimplemented the fold-splitting logic independently. }
\code{eval\_loso.py}\textit{ was extended with a }\code{\_dtw\_fold()}
\textit{ helper that fits the exact production
}\code{DtwRiemannHybridClassifier}\textit{ class (the one
}\code{train\_dtw\_hybrid.py}\textit{ uses) on the same session
splits already used for the flat/hierarchical comparison, so that
all three architectures in Table~\ref{tab:perfold_accuracy} are now
evaluated by one script, on identical folds, rather than being
assembled from two separate evaluation code paths. No change was
made to the training logic, gating, feature extraction, or fold
structure of any architecture --- only reporting was added on top
of the existing evaluation.}

%% ----------------------------------------------------------
\section{Per-Class Performance}
\label{sec:res_perclass}
%% ----------------------------------------------------------

Table~\ref{tab:final_perclass} gives the complete per-class
precision, recall, F1, and support, computed by pooling predictions
across all three \acrshort{loso} folds (\code{eval\_loso.py},
production \code{DtwRiemannHybridClassifier}).

\begin{table}[htbp]
\centering
\caption{Final per-class precision, recall, F1, and support, DTW
hybrid model, LOSO protocol (predictions pooled across all 3
folds). Support totals 1{,}246, matching the activity-gated
dataset size (Chapter~6).}
\label{tab:final_perclass}
\renewcommand{\arraystretch}{1.3}
\begin{tabular}{c l c c c c}
\toprule
\textbf{Cmd} & \textbf{Class} & \textbf{Precision} &
\textbf{Recall} & \textbf{F1} & \textbf{Support} \\
\midrule
O & \code{single\_blink}  & 0.867 & 0.792 & 0.828 & 255 \\
S & \code{double\_blink}  & 0.808 & 0.882 & 0.843 & 262 \\
C & \code{clinch}         & 0.931 & 0.976 & 0.953 & 165 \\
L & \code{bruxism\_left}  & 1.000 & 0.975 & 0.987 & 239 \\
R & \code{bruxism\_right} & 0.988 & 0.975 & 0.981 & 325 \\
\midrule
\multicolumn{2}{l}{\textbf{Accuracy}} & \multicolumn{3}{l}{0.918} &
1{,}246 \\
\multicolumn{2}{l}{\textbf{Macro avg}} & 0.919 & 0.920 & 0.918 &
1{,}246 \\
\multicolumn{2}{l}{\textbf{Weighted avg}} & 0.920 & 0.918 & 0.918 &
1{,}246 \\
\bottomrule
\end{tabular}
\end{table}

\subsection{Reading the Per-Class Pattern}

The per-class scores fall cleanly into two tiers that mirror the
architectural split established in Chapter~9:

\begin{itemize}[leftmargin=*]
    \item \textbf{Muscle classes (C, L, R): F1 0.95--0.99.} These
    are the classes classified via Riemannian covariance geometry
    (Section~9.5). \code{bruxism\_left} reaches perfect precision
    (1.000) --- every window the model calls left bruxism genuinely
    is left bruxism --- with recall of 0.975 accounting for its
    small number of errors. These near-uniformly high scores
    confirm that covariance structure is a robust, cross-session-stable
    representation for the jaw \acrshort{emg} classes.
    \item \textbf{Blink classes (S, O): F1 0.83--0.84.} These are
    the classes classified via \acrshort{dtw} shape-matching
    (Section~9.6). Their precision/recall values are notably
    asymmetric: \code{single\_blink} has high precision (0.867) but
    lower recall (0.792), while \code{double\_blink} has the
    opposite pattern (0.808 precision, 0.882 recall). This is
    exactly what the single-blink contamination finding
    (Chapter~6) predicts: a contaminated \code{single\_blink}
    window that actually contains a multi-blink cluster looks more
    like a \code{double\_blink} to the classifier, which
    systematically pulls \code{single\_blink} windows toward being
    predicted as \code{double\_blink} --- lowering
    \code{single\_blink} recall while inflating \code{double\_blink}
    recall. Section~\ref{sec:res_confusion} confirms this directly
    from the confusion matrix.
\end{itemize}

%% ----------------------------------------------------------
\section{Confusion Patterns}
\label{sec:res_confusion}
%% ----------------------------------------------------------

Table~\ref{tab:confusion_matrix} gives the full confusion matrix,
pooled across all three \acrshort{loso} folds
(\code{eval\_loso.py}, production \code{DtwRiemannHybridClassifier}).
Rows are the true class; columns are the predicted class.

\begin{table}[htbp]
\centering
\caption{Confusion matrix, DTW hybrid model, LOSO protocol,
pooled across all 3 folds (1{,}246 total predictions). Rows: true
class. Columns: predicted class. Diagonal (correct predictions)
in bold.}
\label{tab:confusion_matrix}
\renewcommand{\arraystretch}{1.3}
\begin{tabular}{l c c c c c}
\toprule
\textbf{True $\downarrow$ / Pred $\rightarrow$} &
\code{single\_} & \code{double\_} & \code{clinch} &
\code{bruxism\_} & \code{bruxism\_} \\
 & \code{blink} & \code{blink} & & \code{left} & \code{right} \\
\midrule
\code{single\_blink}  & \textbf{202} & 52  & 1   & 0   & 0   \\
\code{double\_blink}  & 30  & \textbf{231} & 0   & 0   & 1   \\
\code{clinch}         & 0   & 2   & \textbf{161} & 0   & 2   \\
\code{bruxism\_left}  & 1   & 1   & 3   & \textbf{233} & 1   \\
\code{bruxism\_right} & 0   & 0   & 8   & 0   & \textbf{317} \\
\bottomrule
\end{tabular}
\end{table}

\subsection{Reading the Confusion Matrix}

Out of 1{,}246 pooled predictions, 1{,}144 are correct (trace of
the matrix) and 102 are misclassified --- consistent with the
91.8\% pooled accuracy in Section~\ref{sec:res_headline}. The
distribution of those 102 errors confirms every prediction made
qualitatively in earlier chapters:

\begin{itemize}[leftmargin=*]
    \item \textbf{S$\leftrightarrow$O (single/double blink) accounts
    for 82 of the 102 total errors (80.4\%).} 52
    \code{single\_blink} windows are predicted as
    \code{double\_blink}, and 30 \code{double\_blink} windows are
    predicted as \code{single\_blink}. This is by far the dominant
    confusion in the whole matrix and is the direct, quantitative
    signature of the single-blink contamination problem
    (Chapter~6): a \code{single\_blink} recording that actually
    contains a multi-blink cluster produces a waveform shape closer
    to a genuine \code{double\_blink} exemplar, which is exactly
    what \acrshort{dtw} nearest-neighbour classification would be
    expected to do with such an epoch. The asymmetry (52 vs.\ 30)
    is also consistent with Chapter~6's finding that contamination
    specifically afflicts the \code{single\_blink} recordings
    (only 39\% clean) far more than \code{double\_blink}
    (55\% clean) --- more \code{single\_blink} epochs are
    ambiguous, so more \code{single\_blink} epochs are misrouted.
    \item \textbf{Blink$\leftrightarrow$muscle confusion is
    negligible:} only 2 epochs total cross the blink/muscle
    boundary (1 \code{single\_blink}$\rightarrow$\code{clinch}, 1
    \code{bruxism\_left}$\rightarrow$\code{single\_blink}), out of
    1{,}246 predictions --- confirming the Stage~1 router's
    near-perfect isolation (Chapter~9).
    \item \textbf{Confusion among the muscle classes is small and
    concentrated on clinch, not between left and right bruxism.}
    8 \code{bruxism\_right} epochs are misclassified as
    \code{clinch}, and 3 \code{bruxism\_left} epochs as
    \code{clinch}, but \textbf{zero} \code{bruxism\_left} epochs
    are confused with \code{bruxism\_right} or vice versa. This
    matches the architectural expectation exactly: clinch and
    bruxism share the same electrode pair (T3, T4) and are
    distinguished mainly by lateral balance, whereas left and right
    bruxism are separated by the sign of the same asymmetry signal
    and remain essentially unconfusable once routed to the
    Riemannian branch.
\end{itemize}

%% ----------------------------------------------------------
\section{Complete Ablation and Experiment Log}
\label{sec:res_ablation}
%% ----------------------------------------------------------

Table~\ref{tab:experiment_log} consolidates every experiment
referenced across Chapters~6--9 into a single chronological
ablation log, making explicit which changes were kept in the final
production model and which were tried and rejected.

\begin{table}[htbp]
\centering
\caption{Complete experiment log across the project. ``Kept''
experiments are part of the final production model
(\code{dtw\_hybrid\_model.pkl}); ``Rejected'' experiments were
implemented, evaluated, and explicitly not carried forward.}
\label{tab:experiment_log}
\renewcommand{\arraystretch}{1.3}
\begin{tabular}{p{2.1cm} p{4.3cm} p{4.8cm}}
\toprule
\textbf{Phase} & \textbf{Experiment} & \textbf{Outcome} \\
\midrule
Taxonomy & 5$\rightarrow$6 class split (Chapter~6) &
    +3.0\% --- kept \\
Taxonomy & 6$\rightarrow$5 class (drop triple blink, bilateral
    bruxism; Chapter~6) & Cleaner, more robust classes --- kept \\
Epoching & Activity gating, recording-relative (Chapter~7) &
    +3.1\% --- kept \\
Epoching & Session-simulating augmentation
    (jitter/scale/shift/warp) & $-$1.1\% --- \textbf{rejected} \\
Features & \acrshort{dwt}, full (108 feats; Chapter~8) &
    $-$14.4\% --- \textbf{rejected} \\
Features & \acrshort{dwt}, relative-energy only (36 feats;
    Chapter~8) & Break-even --- \textbf{rejected} \\
Features & \acrshort{tkeo} (Chapter~8) & $-$5.7\% ---
    \textbf{rejected} \\
Features & \acrshort{emg} descriptors (Chapter~8) & Hurts flat
    model overall, but helps muscle branch specifically --- kept
    \emph{for Stage~2b only} \\
Architecture & Hierarchical routing, flat$\rightarrow$2-stage
    (Chapter~9) & +5.0\% (GroupKFold) --- kept \\
Muscle & Riemannian Stage~2b (Chapter~9) & +4.1\% (\acrshort{loso})
    --- kept \\
Muscle & \code{tsupdate=True} (Chapter~9) & Best on \acrshort{loso}
    but collapses on GroupKFold --- \textbf{rejected} \\
Blink & Raw peak-count (Chapter~9) & \textbf{Rejected} \\
Blink & $\mu$V peak-count (Chapter~9) & \textbf{Rejected} \\
Blink & Sliding-window count (Chapter~9) & \textbf{Rejected} \\
Blink & Event-locked epoching (Chapter~9) & \textbf{Rejected}
    (contamination diagnosed here) \\
Blink & \acrshort{dtw}-KNN shape-matching (Chapter~9) & +2.1\%
    (\acrshort{loso}) --- \textbf{kept, final} \\
Blink & \acrshort{dtw} barycenter templates (Chapter~9) & 50.5\%
    --- \textbf{rejected} \\
\bottomrule
\end{tabular}
\end{table}

Read as a whole, this log tells a consistent story: every
\textbf{kept} change either fixed a genuine labelling/data-quality
problem (taxonomy, activity gating) or matched a representation to
the physiology of the signal it processes (Riemannian covariance
for amplitude-drift-prone muscle signals, \acrshort{dtw} shape for
contamination-prone blink signals). Every \textbf{rejected} change
either added complexity without addressing the actual failure mode
(more wavelet features, more counting variants) or overfit to one
evaluation protocol at the expense of the other
(\code{tsupdate=True}). This distinction --- between features that
fail because a class is inherently hard, and methods that fail
because they are solving the wrong problem for that class --- is
the throughline of the entire classification methodology developed
in Chapter~9.

%% ----------------------------------------------------------
\section{Real-Time Operating Characteristics}
\label{sec:res_realtime}
%% ----------------------------------------------------------

Although the full real-time inference pipeline is described in
Chapter~12, its headline operating characteristics are reported
here alongside the offline evaluation results for completeness:

\begin{itemize}[leftmargin=*]
    \item \textbf{Per-window latency:} 14\,ms (\acrshort{dtw}
    Stage~2a is the dominant cost), well under the 200\,ms
    real-time step size --- the model comfortably meets real-time
    throughput requirements with substantial margin.
    \item \textbf{Real-time activity gate:} a fixed 20\,\si{\micro
    \volt} \acrshort{rms} threshold at Fp1/Fp2 and T3/T4 (rest
    typically 7--13\,\si{\micro\volt}, active gestures typically
    30--200\,\si{\micro\volt}) --- distinct from the
    recording-relative gate used for offline dataset construction
    (Chapter~7); see Chapter~12 for the reason the two gates
    differ.
    \item \textbf{Confidence threshold:} 0.60, below which no
    command is emitted.
    \item \textbf{\acrshort{fsm} confirmation:} 3 consecutive
    identical predictions required before a command is confirmed
    and transmitted over \acrshort{uart} (Chapter~4).
\end{itemize}

%% ----------------------------------------------------------
\section{Chapter Summary}
\label{sec:res_summary}
%% ----------------------------------------------------------

This chapter presented the detailed results of the final DTW
hybrid model: \textbf{91.3\% accuracy, macro-F1 0.913} as the
unweighted mean across the three \acrshort{loso} folds (Chapter~9),
and \textbf{91.8\% accuracy, macro-F1 0.918} pooled across all
1{,}246 held-out predictions --- both computed under the honest
\acrshort{loso} protocol, differing only in aggregation method.
The per-fold breakdown showed DTW-hybrid outperforming the flat
and hierarchical baselines on every individual session, with
Session3 as the largest and best-performing fold (95.7\%).
Per-class precision/recall/F1 split cleanly into a high-performing
muscle tier (C: 0.953, L: 0.987, R: 0.981, via Riemannian
covariance geometry) and a comparatively harder blink tier (O:
0.828, S: 0.843, via \acrshort{dtw} shape-matching), consistent
with the architecture developed in Chapter~9. The full confusion
matrix confirmed, quantitatively, that 80.4\% of all errors
(82 of 102) are single$\leftrightarrow$double blink confusion ---
the direct signature of the single-blink contamination finding of
Chapter~6 --- while blink$\leftrightarrow$muscle confusion is
negligible (2 epochs total) and left/right bruxism are never
confused with each other. A complete, chronological experiment log
consolidated every kept and rejected change across Chapters~6--9
into a single ablation reference table. Real-time operating
characteristics (14\,ms latency, 0.60 confidence threshold,
3-consecutive \acrshort{fsm} confirmation) were summarised, with
full detail deferred to Chapter~12. The following chapter describes
the configuration and control application used to map classifier
output to prosthetic hand actions.